\documentclass[12pt]{article}
\usepackage[utf8]{inputenc}
\usepackage[T1]{fontenc}
\usepackage{hyperref}
\title{Powari Language Resources for K-12 Curriculum Development}
\author{Compiled by Claude Sonnet 5.5 (claude-sonnet-5-5) for Andr\'as Kornai}
\date{}
\begin{document}
\maketitle

\section{Introduction}
Powari (ISO 639-3 pwr) is an Indo-Aryan language spoken in Balaghat and Seoni districts of Madhya Pradesh and Bhandara and Gondia districts of Maharashtra, written in Devanagari. It has been classified as a variety of Bundeli or as a separate Eastern or Western Hindi language; the Indian census counts it under Hindi, and the Wikipedia figure of about 325,772 speakers (2011) is therefore provisional. Search found only encyclopedic descriptions.

\section{Primary Educational Resources}
None found in search results.

\section{Linguistic and Reference Materials}
\begin{itemize}
\item Wikipedia, Powari language: \url{https://en.wikipedia.org/wiki/Powari_language}
\item Wikipedia, Bagheli language (a related variety on which Powari is said to be based): \url{https://en.wikipedia.org/wiki/Bagheli_language}
\end{itemize}

\section{Religious and Community Resources}
None found.

\section{Digital and Audio Resources}
None found.

\section{Curriculum Development Context}
Powari speakers live among Marathi and Hindi speakers and the language has no known schooling presence. Any curriculum would require fieldwork first: documentation, an agreed Devanagari orthography, and community consultation.

\end{document}